\section{Heat evolution with complete derivative control}
\label{sec:heat}

\subsection{Semigroup and prepared strong solutions}

Functional calculus for the realization in Section~\ref{sec:preparation}
gives
\begin{equation}
 P_n\widehat{\ee^{-\tau L}S}(k)
 =\ee^{-\tau(k^2+4n^2)}P_n\widehat S(k),\qquad
 \norm{\ee^{-\tau L}}=\ee^{-4\tau},\quad \tau\ge0.
 \label{eq:semigroup}
\end{equation}
Dominated convergence proves strong continuity at zero. For $r\ge0$ and
$\tau>0$,
\begin{equation}
 \norm{L^r\ee^{-\tau L}}
 =\sup_{x\ge4}x^r\ee^{-\tau x}<\infty.
 \label{eq:smoothing}
\end{equation}
This gives positive-time left-energy smoothing. It supplies neither
uniform arbitrary-data smoothing at zero nor reference-side regularity.
The semigroup is not compact: even one angular channel contains a
nonzero multiplication operator on a nonatomic radial spectral space.
In particular it is not trace class.

For data in $\Dom L$, the semigroup is the unique strong solution of
$\partial_\tau S=-LS$. Indeed, a difference $U$ of two strong solutions
with zero initial data satisfies
\begin{equation*}
 \frac\dd{\dd\tau}\norm{U(\tau)}^2
 =-2\braket{U(\tau)}{LU(\tau)}
 \le-8\norm{U(\tau)}^2,
\end{equation*}
so its norm remains zero. Define the prepared evolution by
\begin{equation}
 S_{\lambda,\tau}=\ee^{-\tau L}(F_\lambda B).
 \label{eq:evolved-field}
\end{equation}
Proposition~\ref{prop:graph-identity} gives, at every finite order,
\begin{align}
 \partial_\lambda^a\partial_v^h\partial_\tau^cS_{\lambda,\tau}
 &=(-1)^c\ee^{-\tau L}L^c(f_{a,\lambda}^{(h)}B),
 \label{eq:graph-evolution}\\
 \norm{L^p\partial_\lambda^a\partial_v^h\partial_\tau^cS_{\lambda,\tau}}
 &\le\ee^{-4\tau}\norm{L^{p+c}(f_{a,\lambda}^{(h)}B)}.
 \label{eq:graph-evolution-bound}
\end{align}
Strong continuity in the graph spaces, with one extra finite power for
each differentiation, gives the actual right derivatives at zero.
Compactness of the parameterized initial family makes convergence
uniform on $|\lambda|\le M$.

\subsection{Weighted spatial estimates}

Graph control alone does not justify an entire pointwise Fourier
readout. We therefore propagate exponential spatial moments explicitly.
Let
\begin{equation}
 K_\tau(y)=(4\pi\tau)^{-1/2}\ee^{-y^2/(4\tau)},\qquad \tau>0,
 \qquad K_0*f=f.
 \label{eq:heat-kernel}
\end{equation}
The commuting radial and angular multipliers give
\begin{equation}
 S_{\lambda,\tau}(v)
 =(K_\tau*F_\lambda)(v)\,\ee^{-4\tau A}B.
 \label{eq:heat-factorization}
\end{equation}
This is initially an equality in $\HH$. The following estimates also
place both sides in all weighted spaces needed for observation.

The weight $w_{r,b}(v)=(1+|v|)^r\ee^{b|v|}$ is submultiplicative:
$w_{r,b}(x+y)\le w_{r,b}(x)w_{r,b}(y)$. For a standard real Gaussian
$X$, set
\begin{equation}
 \gamma_{r,b}(T)=
 \mathbb E\bigl[(1+\sqrt{2T}|X|)^r\ee^{b\sqrt{2T}|X|}\bigr],
 \qquad \gamma_{r,b}(0)=1.
 \label{eq:gamma}
\end{equation}
Then $\norm{w_{r,b}K_\tau}_1\le\gamma_{r,b}(T)<\infty$ for
$0<\tau\le T$. For an explicit bound, write
$m_\ell=\mathbb E|X|^\ell
=2^{\ell/2}\Gamma((\ell+1)/2)/\sqrt\pi$. Exponential tilting of
the Gaussian gives
\begin{equation}
 \gamma_{r,b}(T)\le2\ee^{b^2T}
 \sum_{j=0}^r\binom rj\sum_{\ell=0}^j\binom j\ell
 (2bT)^{j-\ell}(2T)^{\ell/2}m_\ell.
 \label{eq:gamma-bound}
\end{equation}
To see this, expand $(1+|Y|)^r$ for $Y=\sqrt{2T}X$, bound
$\ee^{b|Y|}$ by $\ee^{bY}+\ee^{-bY}$, and complete the square.
The tilted means are $\pm2bT$, so the triangle and binomial inequalities
give \eqref{eq:gamma-bound}. The convention $0^0=1$ covers the endpoints;
the bound need not be sharp there.

Tonelli's theorem and submultiplicativity now imply
\begin{equation}
 \norm{w_{r,b}(K_\tau*f)}_1
 \le\gamma_{r,b}(T)\norm{w_{r,b}f}_1,
 \qquad 0\le\tau\le T.
 \label{eq:weighted-convolution}
\end{equation}
This estimate propagates every fixed polynomial and exponential
moment. No claim that positive-time convolution preserves the source's
original double-exponential spatial decay is needed.

Move derivatives of the heat kernel onto the smooth, integrable prepared
data. For any finite $a,c,h\ge0$, this yields
\begin{equation}
 \partial_\lambda^a\partial_\tau^c\partial_v^hS_{\lambda,\tau}
 =\sum_{j=0}^c\binom cj(-4)^{c-j}
 (K_\tau*f_{a,\lambda}^{(h+2j)})\,
 A^{c-j}\ee^{-4\tau A}B.
 \label{eq:heat-derivatives}
\end{equation}
Consequently, on the parameter window \eqref{eq:windows},
\begin{align}
 &\norm{w_{r,b}\partial_\lambda^a\partial_\tau^c
                  \partial_v^hS_{\lambda,\tau}}_{L^1(\HS)}
 \le\mathcal M_{a,c,h;r,b}(M,T),
 \label{eq:heat-majorant}\\
 &\mathcal M_{a,c,h;r,b}(M,T)
 :=\gamma_{r,b}(T)\sum_{j=0}^c\binom cj4^{c-j}
       \sqrt{\mu_{2(c-j)}}\,
       \mathcal N_{a,h+2j;r,b}(M)<\infty.
 \label{eq:M-definition}
\end{align}
An additional left power $L^p$ has the same absolute estimate with $c$
replaced by $c+p$. An angular tail replaces the appropriate $\mu$
moment by the vanishing tail in \eqref{eq:bank-tail}. These estimates
justify complete angular and reference limits as well as the radial
integrals.

\subsection{Why the limits at zero are actual derivatives}

For $\tau>0$, integration by parts and the Gaussian kernel justify
\eqref{eq:heat-derivatives}. A bound alone would not establish that
its limit at zero differentiates the original field. The missing step
is continuity in the weighted Banach spaces.

Translation is continuous in $L^1(w_{r,b}\dd v)$. Indeed, approximate
by compactly supported continuous functions; for the remainder, the
translation norm is at most $w_{r,b}(y)$ times its original norm.
Writing $K_\tau*f-f$ as the Gaussian average of translation differences
therefore gives convergence to zero. A dominating random variable is
\begin{equation*}
 \bigl(1+w_{r,b}(\sqrt{2T}X)\bigr)\norm{w_{r,b}f}_1,
\end{equation*}
which is integrable by \eqref{eq:gamma}. This convergence is uniform on
compact sets of $f$ in the weighted space. Section~\ref{sec:source}
established precisely that compactness for
$\{f_{a,\lambda}^{(h)}:|\lambda|\le M\}$.

On the angular side,
$A^s\ee^{-4\tau A}B\to A^sB$ in Hilbert--Schmidt norm by
\eqref{eq:bank-moments} and dominated convergence. Hence every
expression on the right side of \eqref{eq:heat-derivatives} has a
continuous weighted limit at zero. Integrating the positive-time
identity
\begin{equation*}
 \partial_\tau(K_\tau*f)=K_\tau*f''
\end{equation*}
from a positive lower endpoint and then taking that endpoint to zero
shows that its limiting value is the actual right derivative. Iteration
on the controlled derivatives gives every finite order. The same
argument, using the next parameter insertion or spatial derivative,
proves all stated mixed interchanges. Thus
\eqref{eq:heat-majorant} controls continuous integral identities,
including the boundary $\tau=0$.
